\preliminarypart{Notation and Conventions}

General mathematical notation and writing follow \textit{(MATH000) Mathematical Language}. The conventions below record only notation specific to differential equations used in this book. Widely used standard alternatives may be listed in preferred order; the first form is used throughout.

\notationheading{Differential Equations}

The Fourier transform uses the normalization
\[
\widehat f(\boldsymbol{\xi})\coloneqq
\int_{\R^n} f(\mathbf{x})e^{-2\pi i\mathbf{x}\cdot\boldsymbol{\xi}}\,\dd\mathbf{x},
\qquad
f(\mathbf{x})=
\int_{\R^n}\widehat f(\boldsymbol{\xi})e^{2\pi i\mathbf{x}\cdot\boldsymbol{\xi}}\,\dd\boldsymbol{\xi},
\]
whenever the Fourier inversion formula applies.

\begin{notationtable}
$y'$ & First derivative with respect to the understood independent variable. \\
$y''$ & Second derivative. \\
$\dot{x}$ & First derivative of $x$ with respect to time. \\
$\ddot{x}$ & Second derivative of $x$ with respect to time. \\
$y_h,\ y_c$ & Homogeneous or complementary part of a solution. \\
$y_p$ & Particular solution. \\
$c_1,c_2,\dots$ & Arbitrary constants in a general solution. \\
$\begin{aligned} &W(y_1,\dots,y_n)(t),\\ &W[y_1,\dots,y_n](t)\end{aligned}$ & Wronskian of $y_1,\dots,y_n$. \\
$\begin{aligned} &\mathbf{x}'=\mathbf A\mathbf{x}+\mathbf g(t),\\ &\dot{\mathbf{x}}=\mathbf A\mathbf{x}+\mathbf g(t)\end{aligned}$ & Linear first-order system in state-vector form. \\
$\boldsymbol{\Phi}(t)$ & Fundamental matrix of a linear system. \\
\addlinespace[0.26em]
$\begin{aligned} &\mathcal{L}\{f(t)\}(s),\\ &\mathcal{L}[f](s)\end{aligned}$ & Laplace transform of $f$. \\
$\begin{aligned} &\mathcal{L}^{-1}\{F(s)\}(t),\\ &\mathcal{L}^{-1}[F](t)\end{aligned}$ & Inverse Laplace transform of $F$. \\
$(f\ast g)(t)$ & One-sided convolution, $\displaystyle (f\ast g)(t)\coloneqq\int_0^t f(\tau)g(t-\tau)\,\dd\tau$. \\
$H(t-a),\ u(t-a)$ & Heaviside or unit-step function shifted to $t=a$; its value at the jump is specified when relevant. \\
\addlinespace[0.26em]
$G(x,\xi)$ & Green's function; the variables depend on context. \\
$J_\nu(x)$ & Bessel function of the first kind of order $\nu$. \\
$Y_\nu(x),\ N_\nu(x)$ & Bessel function of the second kind (Neumann function) of order $\nu$. \\
\addlinespace[0.26em]
$\widehat f(\boldsymbol{\xi}),\ \mathcal{F}f(\boldsymbol{\xi})$ & Fourier transform under the normalization above. \\
\end{notationtable}
